\documentclass[conference]{IEEEtran}
\IEEEoverridecommandlockouts
% The preceding line is only needed to identify funding in the first footnote. If that is unneeded, please comment it out.
\usepackage{cite}
\usepackage{amsmath,amssymb,amsfonts}
\usepackage{algorithmic}
\usepackage{graphicx}
\usepackage{textcomp}
\usepackage{xcolor}
\def\BibTeX{{\rm B\kern-.05em{\sc i\kern-.025em b}\kern-.08em
    T\kern-.1667em\lower.7ex\hbox{E}\kern-.125emX}}
\begin{document}

\title{Conference Paper Title*\\
{\footnotesize \textsuperscript{*}Note: Sub-titles are not captured in Xplore and
should not be used}
\thanks{Identify applicable funding agency here. If none, delete this.}
}

\author{\IEEEauthorblockN{1\textsuperscript{st} José Raimundo Alves dos Santos Júnior}
\IEEEauthorblockA{\textit{dept. name of organization (of Aff.)} \\
\textit{name of organization (of Aff.)}\\
City, Country \\
email address or ORCID}
\and
\IEEEauthorblockN{2\textsuperscript{nd} Eduardo Furtado de Simas Filho}
\IEEEauthorblockA{\textit{dept. name of organization (of Aff.)} \\
\textit{name of organization (of Aff.)}\\
City, Country \\
email address or ORCID}

}

\maketitle

\begin{abstract}
This document is a model and instructions for \LaTeX.
This and the IEEEtran.cls file define the components of your paper [title, text, heads, etc.]. *CRITICAL: Do Not Use Symbols, Special Characters, Footnotes, 
or Math in Paper Title or Abstract.
\end{abstract}

\begin{IEEEkeywords}
component, formatting, style, styling, insert
\end{IEEEkeywords}

\section{Introduction}
O monitoramento de progresso é um dos elos fundamentais para gestão de projetos, auxiliando na conclusão bem sucedida da construção de edificações e infraestruturas (Pal et al., 2023). Na construção de Obras de Arte Especiais (OAE), como as pontes, as condições diferem significativamente das edificações convencionais, caracterizando-se por processos altamente setorizados, distribuição espacial discreta e uso intensivo de componentes pré-fabricados (Zhang, 2023). Embora essas estruturas apresentem padrões geométricos repetitivos e similares, a extensão do canteiro e a complexidade logística impõem desafios severos à gestão e ao acompanhamento físico do projeto (Zhang, 2023).   Os métodos tradicionais de monitoramento do progresso baseiam-se em registros manuais, inspeções in loco e atualizações em planilhas ou papel (Pal et al., 2023). Essa abordagem apresenta desvantagens críticas: é excessivamente morosa, trabalhosa, propensa a erros humanos e influenciada pela subjetividade dos inspetores (Qureshi et al., 2023). Devido à incapacidade de abranger a totalidade do canteiro com alta frequência, os dados coletados ficam rapidamente desatualizados, tornando a tomada de decisão reativa e propiciando atrasos e estouros de orçamento (Reja et al., 2022). Ademais, a simples captura visual do canteiro sem processamento automatizado não soluciona o problema, pois apenas transfere o gargalo operacional do campo para o escritório.   Nesse cenário, a integração de ferramentas como a Modelagem da Informação da Construção (BIM) a técnicas de Visão Computacional e Inteligência Artificial tem se consolidado como a via mais eficiente para o monitoramento automatizado (Pal et al., 2023). O uso de Aprendizado Profundo (Deep Learning) permite processar dados visuais complexos, como nuvens de pontos, comparando o modelo planejado (as-planned) com a execução real (as-built) para extrair informações objetivas de avanço (Reja et al., 2022; Braun et al., 2020).   Apesar desses avanços, a literatura e as aplicações de monitoramento automatizado ainda concentram-se majoritariamente em edificações verticais, revelando uma lacuna no desenvolvimento de fluxos adaptados às especificidades geométricas e operacionais das pontes. Portanto, este trabalho propõe um fluxo semi-automatizado para o monitoramento de pontes, integrando fotogrametria, segmentação de nuvens de pontos por IA (RandLA-Net) e algoritmos de alinhamento geométrico (OBB-ICP) ao modelo BIM 4D, visando eliminar a carga manual, fornecer dados acionáveis aos gestores e mitigar ineficiências construtivas.   


\section*{References}

Please number citations consecutively within brackets \cite{b1}. The 
sentence punctuation follows the bracket \cite{b2}. Refer simply to the reference 
number, as in \cite{b3}---do not use ``Ref. \cite{b3}'' or ``reference \cite{b3}'' except at 
the beginning of a sentence: ``Reference \cite{b3} was the first $\ldots$''

Number footnotes separately in superscripts. Place the actual footnote at 
the bottom of the column in which it was cited. Do not put footnotes in the 
abstract or reference list. Use letters for table footnotes.

Unless there are six authors or more give all authors' names; do not use 
``et al.''. Papers that have not been published, even if they have been 
submitted for publication, should be cited as ``unpublished'' \cite{b4}. Papers 
that have been accepted for publication should be cited as ``in press'' \cite{b5}. 
Capitalize only the first word in a paper title, except for proper nouns and 
element symbols.

For papers published in translation journals, please give the English 
citation first, followed by the original foreign-language citation \cite{b6}.

\begin{thebibliography}{00}
\bibitem{b1} G. Eason, B. Noble, and I. N. Sneddon, ``On certain integrals of Lipschitz-Hankel type involving products of Bessel functions,'' Phil. Trans. Roy. Soc. London, vol. A247, pp. 529--551, April 1955.
\bibitem{b2} J. Clerk Maxwell, A Treatise on Electricity and Magnetism, 3rd ed., vol. 2. Oxford: Clarendon, 1892, pp.68--73.
\bibitem{b3} I. S. Jacobs and C. P. Bean, ``Fine particles, thin films and exchange anisotropy,'' in Magnetism, vol. III, G. T. Rado and H. Suhl, Eds. New York: Academic, 1963, pp. 271--350.
\bibitem{b4} K. Elissa, ``Title of paper if known,'' unpublished.
\bibitem{b5} R. Nicole, ``Title of paper with only first word capitalized,'' J. Name Stand. Abbrev., in press.
\bibitem{b6} Y. Yorozu, M. Hirano, K. Oka, and Y. Tagawa, ``Electron spectroscopy studies on magneto-optical media and plastic substrate interface,'' IEEE Transl. J. Magn. Japan, vol. 2, pp. 740--741, August 1987 [Digests 9th Annual Conf. Magnetics Japan, p. 301, 1982].
\bibitem{b7} M. Young, The Technical Writer's Handbook. Mill Valley, CA: University Science, 1989.
\end{thebibliography}
\vspace{12pt}
\color{red}
IEEE conference templates contain guidance text for composing and formatting conference papers. Please ensure that all template text is removed from your conference paper prior to submission to the conference. Failure to remove the template text from your paper may result in your paper not being published.

\end{document}

